\documentclass{chinese-journal-of-mechanical-engineering}

%%%%%%%%%%%%%%%%%%%%%%%%%%%%%%%%%%%%%%%%%%%%%%%%%%%%%%%%%%%%%
% These packages below are only for test.
% You can remove them in real paper writing after removing
% the test commands in the main text.
%%%%%%%%%%%%%%%%%%%%%%%%%%%%%%%%%%%%%%%%%%%%%%%%%%%%%%%%%%%%%
\usepackage{lipsum}
\usepackage{mwe}
\usepackage{zhlipsum}

%%%%%%%%%%%%%%%%%%%%%%%%%%%%%%%%%%%%%%%%%%%%%%%%%%%%%%%%%%%%%
% Here is your bib.
%%%%%%%%%%%%%%%%%%%%%%%%%%%%%%%%%%%%%%%%%%%%%%%%%%%%%%%%%%%%%
\addbibresource{main.bib}

%%%%%%%%%%%%%%%%%%%%%%%%%%%%%%%%%%%%%%%%%%%%%%%%%%%%%%%%%%%%%
% Input the DOI and other stuff.
%%%%%%%%%%%%%%%%%%%%%%%%%%%%%%%%%%%%%%%%%%%%%%%%%%%%%%%%%%%%%
\doi{10.3901/JME.20**.**.***}

\renewcommand{\InfoBundle}{%
    \title{这是论文题目}%
    \author[1]{作者一}%
    \author[2]{作者二}%
    \author[1,2]{作者三}%
    \affil[1]{机构一}%
    \affil[2]{机构二}%
    \abst{\zhlipsum[1]}%
    \keyword{关键词一；关键词二；关键词三}}

\renewcommand{\EnInfoBundle}{%
    \title{This Is The English Title}%
    \author[1]{Author I}%
    \author[2]{Author II}%
    \author[1,2]{Author III}%
    \affil[1]{Institute A}%
    \affil[2]{Institute B}%
    \abst{\lipsum[1]}%
    \keyword{Keyword 1; Keyword 2; Keyword 3}}

%%%%%%%%%%%%%%%%%%%%%%%%%%%%%%%%%%%%%%%%%%%%%%%%%%%%%%%%%%%%%
% Below is the main document.
%%%%%%%%%%%%%%%%%%%%%%%%%%%%%%%%%%%%%%%%%%%%%%%%%%%%%%%%%%%%%    
\begin{document}

% Title, abstract, keywords, etc.
\maketitle
\thispagestyle{firststyle}

% Main text
\setcounter{section}{-1}
\begin{multicols}{2}

    %%%%%%%%%%%%%%%%%%%%%%%%%%%%%%%%%%%%%%%%%%%%%%%%%%%%%%%%%%%%%
    \section{前言}
    \zhlipsum[1]

    %%%%%%%%%%%%%%%%%%%%%%%%%%%%%%%%%%%%%%%%%%%%%%%%%%%%%%%%%%%%%
    \section{论文常见元素测试}
    以下是多段落纯文本测试。

    \zhlipsum[1-3]

    以下是次级标题测试。

    \subsection{引用参考文献测试}
    这是引用\cite{aikenhead1983canadarm}。

    \subsection{有序列表测试}
    这是有序列表测试。

    \begin{enumerate}[（1）]
        \item \zhlipsum[1]
        \item 第二个项目。
        \item 第三个项目。
    \end{enumerate}

    \subsection{图片测试}
    以下是图片测试，见图 \ref{fig:example-1}。

    \begin{figure}[H]
        \centering
        \includegraphics[width=\columnwidth]{example-image-a}
        \caption{测试图片}\label{fig:example-1}
    \end{figure}

    以下是占据两列的图片测试，见图 \ref{fig:example-2}，为了更好地展示效果，将插入一些示例文本。

    \begin{figure*}
        \centering
        \includegraphics[width=.5\textwidth]{example-image-a}
        \caption{测试多列图片}\label{fig:example-2}
    \end{figure*}

    特别注意，这里占据两列的图片使用了 \texttt{figure*} 环境。关于这一环境产生的浮动体的位置问题，参见 \url{https://tex.stackexchange.com/a/167203}。使用类似的 \texttt{table*} 同样可以生成多列表格。

    \zhlipsum[1]

    \subsection{表格测试}
    以下是表格测试，见表 \ref{tab:example-1}。

    \begin{table}[H]
        \caption{测试表格}\label{tab:example-1}
        \centering
        \begin{tabular}{ccc}\toprule
            表头 1 & 表头 2 & 表头 3 \\\midrule
            1 & 2 & 3 \\
            4 & 5 & 6 \\\bottomrule
        \end{tabular}
    \end{table}

    \subsection{公式测试}
    以下是公式测试，见公式 (\ref{eq:example-1}) 和 (\ref{eq:example-2})。
    \begin{gather}
        \bm{F}=m\frac{\mathrm{d}\bm{r}}{\mathrm{d}t} \label{eq:example-1} \\
        \begin{cases}
            a+b=2 \\
            2a+b=4
        \end{cases} \label{eq:example-2}
    \end{gather}

    在 \texttt{multicols} 多列环境中尚未找到生成多列公式的较好方法，除非退出 \texttt{multicols} 环境后插入公式。这里不作测试。

    \subsection{算法测试}
    以下是算法测试，见算法 \ref{alg:example-1}。

    特别注意，这里使用了 \texttt{algorithm2e} 宏包，可能和规范有所不符。如有需要可以自行实现。

    \begin{algorithm}[H]
        \caption{测试算法}\label{alg:example-1}
        \KwIn{输入}
        \KwOut{输出}
        步骤 1\;
        步骤 2\;
        步骤 3\;
        步骤 4\;
    \end{algorithm}

    %%%%%%%%%%%%%%%%%%%%%%%%%%%%%%%%%%%%%%%%%%%%%%%%%%%%%%%%%%%%%
    \section{结论}
    这里为了更好看，填充一些文本。

    \zhlipsum[3]
    
    %%%%%%%%%%%%%%%%%%%%%%%%%%%%%%%%%%%%%%%%%%%%%%%%%%%%%%%%%%%%%
    \printbibliography

    %%%%%%%%%%%%%%%%%%%%%%%%%%%%%%%%%%%%%%%%%%%%%%%%%%%%%%%%%%%%%
    % Author information here
    \hrule\setlength{\parindent}{0pt}\vskip 1em
    \begingroup\footnotesize
    {\sffamily 作者简介：} % Don't leave a line break here
    作者一，男，19xx年出生，博士研究生。 

    E-mail：\href{mailto:author-i@example.com}{author-i@example.com}

    作者二，女，19xx年出生，博士，教授，博士研究生导师。

    E-mail：\href{mailto:author-ii@example.com}{author-ii@example.com}

    作者三，...。

    E-mail：\href{mailto:author-iii@example.com}{author-iii@example.com}

    % If you fail to keep the same vertical gaps, try to add this line below.
    % \phantom{}
    \endgroup

\end{multicols}

\end{document}
